\section{Behavioral analysis with executable conditions}
\label{sec:executable-analysis}
The development studies augment retrieved guidance with task-specific
programs to examine whether a repair preserves the intended behavior.
Figures~\ref{fig:concept} and~\ref{fig:behavior-workflow} show their comparison and workflow
in the main text; this appendix gives the construction protocol.
\paragraph{Values alone do not describe a condition.}
In Python, a named tuple is a tuple whose fields also have names. In our
examples, constructing it takes two arguments; an ordinary tuple subclass
takes one iterable. Both can hold the same values but exercise
different code. Likewise, a timezone can describe the stored data or the
requested output. These checks depend on the input roles as well as their values.

\paragraph{Two passing examples can hide a failing combination.}
Our Django case uses a lazy User object: a wrapper that loads the user when
needed. A query using this object passes, as does an older range query using
a named tuple to hold its endpoints. Putting the lazy object inside the
named-tuple range in the current query fails. This motivates
\emph{condition composition}: inserting a historical trigger into the current operation
(Figure~\ref{fig:concept}).

\subsection{Executable-condition study protocol}
\label{sec:executable-protocol}
The executable-analysis extension connects example construction to a repair loop
(Figure~\ref{fig:behavior-workflow}). Construction records what a past repair taught
and which conditions make that lesson relevant. During repair, retrieval
selects useful experience, composition brings its conditions into the
current workflow, and execution turns them into feedback for the next edit.
The graph preserves the links between these stages, so the agent can trace
a retrieved requirement to both its supporting example and its expected behavior.

\subsection{Retain the context that makes advice useful}\label{sec:store}
For these studies, an analysis record augments a retrieved requirement with
a test example and its results before and after the historical repair:
\begin{equation}
 m=(r,W,E^{-},E^{+}),\qquad W=(I,A,O),
 \label{eq:witness}
\end{equation}
where $r$ states the requirement, $I$ constructs its triggering input, $A$
performs the operation, and $O$ observes the expected outcome. We call $W$
a witness: an executable example that fails on the historical code $V^{-}$
and passes on its repaired version $V^{+}$. $E^{-}$ and $E^{+}$ record these executions.

A language-model writer extracts the requirement and example from the
historical issue, patch, and relevant code. The example is checked against
the historical repair, then stored with its requirement and execution results.
The graph also records how inputs participate in an operation: which values
form a range, for example, or which timezone specifies the output. These
relations preserve the information needed to adapt an experience after retrieval.

\paragraph{Retrieve the requirement together with its context.}
The study retrieves relevant historical repairs, then follows record links
to return their requirements, programs, execution results, and input roles
together. A useful source may concern a related operation: a fix for deleting
one kind of database constraint can help repair deletion of another kind.
Keeping the requirement connected to its context gives the agent both a
lesson and the information needed to apply it. These studies use
similarity search; Appendix~\ref{sec:method-contract} gives the retrieval
settings.

\subsection{Check a historical requirement in the current workflow}
\label{sec:compose}
Composition turns retrieved experience into a check for the current task.
It starts with $P_q$, a runnable program that reproduces the current issue,
and a historical requirement applicable to the current code. An
\emph{adapter}, a small program-transformation routine, uses the stored input
relations to insert the historical trigger into $P_q$. The check keeps the
current workflow and tests the expected behavior under this added condition.

For example, a historical range type can be introduced into the current
query while preserving its input objects and query structure
(Figure~\ref{fig:concept}). Running this combination exposes interactions
that neither the original query nor the historical example exercises alone.
The same design applies to timezone conversion: the historical destination
requirement is checked together with the current source-timezone behavior.
Appendix~\ref{sec:method-contract} gives the input preparation and adapter
constructions.

\subsection{Use failure observations to guide the next edit}\label{sec:feedback}
The agent runs the constructed checks and uses failures to revise the patch.
Each failure identifies an input and expected behavior that the repair needs
to satisfy. Rerunning the same check gives successive edits a concrete
repair objective.
